\initchapter{Estado del arte}\label{chap:state art}

En este capítulo se examinan a los trabajos desarrollados hasta el momento sobre la aplicación de controladores en vehículos aeroespaciales, donde los esquemas de control integran a las interfaces cerebro-computadora (\glspl{bci}), la teoría \textit{wavelet} y las \glspl{rna} para el análisis de las señales de control y de \acrfull{eeg}.

Se inicia mencionado a las áreas de interés para el desarrollo de esta tesis en la \cref{sec:art intro} llamada \textit{Introducción}. Después, en el apartado \ref{sec:art relevant research} titulado \textit{Trabajos relevantes}, se expone parte de las publicaciones científicas consultadas haciendo énfasis en sus áreas de aplicación y técnicas utilizadas. Finalmente, en \textit{Conclusión} (v. sec. \ref{sec:art comments}) se ofrecen algunas observaciones que definen la línea de salida para este proyecto.

\section{Introducción}\label{sec:art intro}
    
    En el presente apartado se examinan a los trabajos que hablan sobre la implementación de las interfaces cerebro-computadora en esquemas de control para vehículos aeroespaciales. Buscando que los controladores sean del tipo \acrfull{pmr} e ``inteligentes'', es decir, que integren a las redes neuronales en su estructura interna \cite[p. 285]{Matilde2011} y que, además, utilicen al \acrfull{mra} para el tratamiento de las señales de \acrlong{eeg}. 
    
    Primero se revisarán los casos de estudio hallados sobre aplicaciones en el sector aeroespacial, esto con el fin de evidenciar la ausencia de información relativa al uso de los sistemas \acrshort{bci} en el modelo del helicóptero \textit{Quanser}. Posteriormente, se hablarán de algunos trabajos desarrollados en el campo médico.
    
\section{Trabajos relevantes}\label{sec:art relevant research}

    A continuación se exponen las reseñas de algunos trabajos obtenidos en repositorios. Esta bibliografía consta de tesis de grado, de maestría y de doctorado; así como artículos de revistas y conferencias.
    
    \subsection*{Aplicaciones en el sector aeroespacial}
    
    En el \citeyear{Duan2019}, \citet{Duan2019} desarrollaron en China una \acrlong{bci} para el control de un cuadricóptero usando los \acrfull{mi} y el \acrfull{ssvep}. Una mujer y ocho varones participaron en los experimentos, cuyo promedio de edades era de 25 años. Se utilizó una red de sensores electroencefalográficos \textit{Geodesic} y a la plataforma BCI2000\footnote{Plataforma de código abierto creada por el \acrlong{ncan} (\acrshort{ncan}\textregistered, por sus siglas en inglés).} para la recolección de los impulsos eléctricos de las neuronas, decodificándolas en cinco instrucciones de vuelo: derecha, izquierda, arriba, abajo y \textit{hover}. Los usuarios demostraron su capacidad para manipular mentalmente al vehículo en el mundo real; asimismo, el autor indica que manejar una interfaz híbrida incrementa la cantidad de comandos disponibles y se mejora la tasa de identificación de los mismos. Sin embargo, se excluye el uso del análisis multiresolución dado que se empleó al \acrfull{cca} para el tratamiento de las señales.
    
    Durante el \citeyear{MamaniMusaja2018}, en Perú se creó una \acrlong{bci} para el control de vuelo del cuadricóptero \textit{Bebop} de la compañía Parrot\textregistered\,y \citet{MamaniMusaja2018} optó por utilizar el dispositivo electroencefalográfico \textit{Insight} de la empresa Emotiv. Se aprovecharon a los entornos de desarrollo de cada fabricante así como algunas librerías implementadas en el lenguaje de programación C++ para diseñar de las vistas del programa y del algoritmo de tratamiento de señales. Adicionalmente, se elaboró una aplicación móvil con Andriod Studio\textregistered. En las pruebas participaron cinco personas, aunque se desconocen sus edades y su género. Tampoco se menciona el paradigma utilizado y, pese a que se habla de ``comandos mentales'', el usuario debe mover su cabeza para que el equipo interprete sus órdenes. No obstante, se validaron siete instrucciones: \textit{hover}, adelante, atrás, izquierda, derecha, arriba y abajo. Finalmente, con base en los resultados mostrados se determinó que el sistema funciona con un 30 \% de efectividad, de ahí que el autor recomiende agregar técnicas de inteligencia artificial a su trabajo.

    Por otro lado, en este mismo año \citet{Nourmohammadi2018} presentaron un sondeo comparativo de varias investigaciones realizadas entre 2010 y 2015. De acuerdo con la información revisada, en la mayoría de los casos se utilizaron cuadricópteros (v. \cite{Khan2015,Kim2014a,LaFleur2013}), algunos optaron por helicópteros simulados (v. \cite{Doud2011} y \cite{Royer2010}), hexacópteros \cite{Shi2015} o aeronaves de ala fija \cite{Akce2010}. Todos ellos se basan en el uso de las señales de \acrlong{eeg} generadas por medio de los \acrlong{mi}, el parpadeo del ojo y/o la gesticulación facial; analizándolas a través de clasificadores lineales o binarios y, en otros cosos, con técnicas de \acrlong{ia}, tales como \glspl{rnn} \cite{KosMyna2015} y \glspl{svm} \cite{Kim2014a}.
    
    Por último, cabe mencionar que hay investigaciones donde se aprovechan a las \glspl{wnt} en esquemas de identificación y control para distintos sistemas (v. \cite{CarrilloSantos2012,HernandezPerez2018,VegaNavarrete2013,DominguezMayorga2017,GarciaCastro2020}).
    
    \subsection*{Aplicaciones en el sector médico}

    En el \citeyear{OlivaresCarillo2017}, se expuso un trabajo nacional enfocado en el diseño e implementación de un prototipo para controlar una silla de ruedas. En él, \citet{OlivaresCarillo2017} presenta una serie de protocolos para la adquisición y procesamiento de señales de \acrlong{eeg}, \acrfull{ecg} y \acrfull{eog} generadas por el \acrlong{ssvep}. Se recopiló la actividad cerebral de 14 personas por medio del equipo \textit{Biosemi} y se crearon un par de bancos de información que se trataron con las técnicas del \acrfull{lda} y \acrlong{svm}, esta última fue seleccionada para la aplicación final. Se adaptó un microcontrolador a una silla de ruedas ya motorizada, pues la idea era sustituir los comandos de control propios del aparato por las instrucciones del sistema \acrshort{bci}.
    
    Al mismo tiempo, en España se presentó un controlador para un entorno hospitalario, permitiéndole a un paciente con discapacidad motriz realizar tareas ``cotidianas'', concreta- mente: encender y apagar luces, televisores y la climatización de la habitación, además de modificar la posición de la camilla y las persianas. \citet{SanchezGarcia2017} logró esto manejando a las señales de \acrlong{eeg} como instrucciones de operación para una serie de actuadores programados en el \acrfull{labview}. Se utilizaron ochos electrodos instalados en el equipo \textit{Enobio} para recopilar la información cerebral y se analizó en la plataforma BCI2000. El autor complementa su trabajo dando un presupuesto sobre este proyecto; el cual incluye el software, el hardware y los salarios de las personas involucradas, sumando un total aproximado de \euro 24,896.52\footnote{Valor ajustado de acuerdo con la tasa de inflación consultada el 27/12/2021 en \url{https://www.dineroeneltiempo.com}. La cantidad original dada en el 2017 fue de \euro 23,973.41 \cite[Tab. 11]{SanchezGarcia2017}.}.
    
    Para terminar, en el \citeyear{GarciaBlancas2016} \citet{GarciaBlancas2016} propone la implementación de un sistema \acrshort{bci} para la manipulación de un brazo robótico de \acrlong{2gdl}, aprovechando el análisis multiresolución para el estudio de las señales de \acrshort{eeg} y para el controlador \acrlong{pmr}. La actividad neuroeléctrica generada por los \acrlong{mi} se recopiló a través del equipo portátil Emotiv EPOC+ y el tratamiento consistió en la pre-amplificación, amplificación y filtrado de cada señal, permitiendo la extracción de sus frecuencias características por medio de la \acrfull{dwt}, lo que da lugar a la interpretación de las instrucciones dadas por el usuario. Con respecto al control \acrshort{pmr}, se manejó un esquema de identificación basado en redes neuronales de base radial (\acrshort{rbf}) con funciones de activación \textit{wavelet}, agregando una estructura de filtros de \acrfull{iir} para facilitar la aplicación de un algoritmo de auto-sintonización. Por último, el autor encontró que, gracias a las técnicas aplicadas, basta con utilizar un solo electrodo del dispositivo electroencefalográfico para ejecutar los comandos mentales y, a la vez, apunta al uso de otra red neuronal para la caracterización de las señales cerebrales.
    
    % \matrixstate
    %     {OlivaresCarillo2017} % Cita
    %     {Diseñar e implementar el prototipo funcional para el control de una silla de ruedas usando señales de \acrshort{eeg}} % Objetivos
    %     {Seis y ocho sujetos de prueba para cada banco de datos} % Muestra
    %     {Se realizaron dos bancos de datos con seis y ocho personas, respectivamente. Las señales se filtraron por medio un filtro pasa-altas de 1 Hz. Se evaluó el desempeño de LDA y SVM, se elije al segundo método} % Metodología
    %     {Equipo Biosemi, Se utilizó \acrshort{ssvep}, microcontrolador hacia la silla de ruedas, monitor a 90 cm de distancia} % Técnica
    %     { } % Resultados
    %     { } % Referencias 
    
\section{Conclusión}\label{sec:art comments}
    
    Con base en la literatura encontrada, es posible verificar que hay una carencia de información sobre la aplicación de interfaces cerebro-computadora en el modelo del helicóptero \textit{Quanser} ya que suelen utilizarse otros modelos de aeronaves.
    
    De igual manera, solo se encontró un trabajo que incorpore a los sistemas \acrshort{bci}, el \acrlong{mra} y los controladores proporcionales multiresolución (v. \cite{GarciaBlancas2016}). No obstante, de manera adicional se utilizarán los resultados mostrados por \cite{MamaniMusaja2018,HernandezPerez2018,OlivaresCarillo2017} y \cite{Goeksu2018} para complementar la información sobre las interfaces cerebro-computadora y también se recurrirá a la documentación necesaria para el estudio de la fisiología del cerebro humano.
    
    Referente a los esquemas de control, se tomará la información dada en \cite{VegaNavarrete2013,CarrilloSantos2012} y \cite{VillarealGrajales2017} para la implementación del controlador proporcional multiresolución inteligente en el modelo a trabajar.
    
    Finalmente, se recuerda que esta tesis es una continuación del trabajo presentado en el \citeyear{GarciaCastro2020} titulado \citetitle{GarciaCastro2020} \cite{GarciaCastro2020}, siendo esta la línea de partida.